%!TEX root = ../main.tex
% Chapter 7: Transformers
\chapter{Transformers}
\label{ch:transformers}

\paragraph{Chapter objectives}
\begin{itemize}
\item Understand the Transformer architecture and the attention mechanism
\item Derive the formulas for multi-head attention and positional encoding
\item Distinguish encoder, decoder, and encoder-decoder variants
\item Implement a complete Transformer in PyTorch
\item Know the foundational models: BERT, GPT, ViT
\end{itemize}

\section{``Attention Is All You Need'' --- The paradigm shift}
\label{sec:transformers-intro}

In 2017, Vaswani et al.\ published ``Attention Is All You Need''\index{Attention Is All You Need}, a paper that radically transformed the landscape of natural language processing and, more broadly, deep learning\index{Transformer}. The central idea is to entirely replace recurrences (RNN, LSTM, GRU --- cf.\ previous chapter) with a pure \emph{attention}\index{attention} mechanism, enabling massive parallelism and efficient modeling of long-range dependencies.

\begin{remark}
Before Transformers, sequence-to-sequence\index{sequence-to-sequence} models relied on RNN encoder-decoders with attention (Bahdanau, 2014). The Transformer eliminates all recurrence, retaining only attention.
\end{remark}

\section{Attention mechanism}
\label{sec:attention-mechanism}

\subsection{Scaled Dot-Product Attention}

Let an input sequence of length $n$ with vectors of dimension $d_k$. Three linear projections are defined:
\begin{itemize}
    \item \textbf{Query}: $Q = XW^Q$, with $W^Q \in \R^{d_{\text{model}} \times d_k}$
    \item \textbf{Key}: $K = XW^K$, with $W^K \in \R^{d_{\text{model}} \times d_k}$
    \item \textbf{Value}: $V = XW^V$, with $W^V \in \R^{d_{\text{model}} \times d_v}$
\end{itemize}

\begin{mathbox}[title=\textbf{Scaled Dot-Product Attention}]\index{scaled dot-product attention}
\begin{equation}
\text{Attention}(Q, K, V) = \softmax\!\left(\frac{QK^\top}{\sqrt{d_k}}\right) V
\label{eq:scaled-attention-transformer}
\end{equation}
\end{mathbox}

\begin{intuition}
The factor $\frac{1}{\sqrt{d_k}}$ prevents dot products from becoming too large in high dimensions, which would push the $\softmax$ into regions of near-zero gradient.

Indeed, if the components of $Q$ and $K$ are i.i.d.\ with mean 0 and variance 1, then $\E[q_i^\top k_j] = 0$ and $\Var(q_i^\top k_j) = d_k$. Dividing by $\sqrt{d_k}$ brings the variance back to 1.
\end{intuition}

\subsection{Multi-head attention}
\label{sec:multi-head-attention}

Rather than a single attention function with $d_{\text{model}}$ dimensions, $h$ parallel ``heads''\index{multi-head attention} are used:

\begin{keyformulas}[title=\textbf{Multi-head attention}]
\begin{align}
\text{MultiHead}(Q, K, V) &= \text{Concat}(\text{head}_1, \ldots, \text{head}_h)\, W^O \\
\text{where} \quad \text{head}_i &= \text{Attention}(QW_i^Q,\; KW_i^K,\; VW_i^V)
\end{align}
with $W_i^Q \in \R^{d_{\text{model}} \times d_k}$, $W_i^K \in \R^{d_{\text{model}} \times d_k}$, $W_i^V \in \R^{d_{\text{model}} \times d_v}$, $W^O \in \R^{hd_v \times d_{\text{model}}}$, and typically $d_k = d_v = d_{\text{model}} / h$.
\end{keyformulas}

\begin{remark}
Each head can learn to focus on a different type of relationship: some heads capture syntax, others semantics, and others positional relationships.
\end{remark}

\section{Positional encoding}
\label{sec:positional-encoding}

Since the Transformer has no recurrence, it has no intrinsic notion of token order. A \emph{positional encoding}\index{positional encoding} is therefore added to the input embeddings.

\subsection{Derivation of the sinusoidal formula}

\begin{definition}[Sinusoidal positional encoding]
For a position $\text{pos}$ and a dimension $i$:
\begin{align}
PE_{(\text{pos}, 2i)} &= \sin\!\left(\frac{\text{pos}}{10000^{2i/d_{\text{model}}}}\right) \label{eq:pe-sin}\\
PE_{(\text{pos}, 2i+1)} &= \cos\!\left(\frac{\text{pos}}{10000^{2i/d_{\text{model}}}}\right) \label{eq:pe-cos}
\end{align}
\end{definition}

\begin{theorem}[Relative position property]
For any fixed offset $k$, there exists a linear transformation $M_k$ independent of $\text{pos}$ such that:
\begin{equation}
PE_{\text{pos}+k} = M_k \cdot PE_{\text{pos}}
\end{equation}
\end{theorem}

\begin{proof}
Consider the dimension pair $(2i, 2i+1)$ and let $\omega_i = 1/10000^{2i/d_{\text{model}}}$. We have:
\begin{align}
\begin{pmatrix} \sin(\omega_i(\text{pos}+k)) \\ \cos(\omega_i(\text{pos}+k)) \end{pmatrix}
&= \begin{pmatrix} \cos(\omega_i k) & \sin(\omega_i k) \\ -\sin(\omega_i k) & \cos(\omega_i k) \end{pmatrix}
\begin{pmatrix} \sin(\omega_i \cdot \text{pos}) \\ \cos(\omega_i \cdot \text{pos}) \end{pmatrix}
\end{align}
by the trigonometric addition formulas. The rotation matrix depends only on $k$, not on $\text{pos}$, which allows the model to learn to exploit relative positions.
\end{proof}

\section{Complete Transformer architecture}
\label{sec:transformer-architecture}

\subsection{Encoder block}
\label{sec:encoder-block}

Each encoder block consists of:
\begin{enumerate}
    \item \textbf{Multi-Head Self-Attention}
    \item \textbf{Add \& Norm} (residual connection + layer normalization)
    \item \textbf{Feed-Forward Network} (FFN) with two layers:
    \begin{equation}
    \text{FFN}(x) = \relu(xW_1 + b_1)W_2 + b_2
    \end{equation}
    with $W_1 \in \R^{d_{\text{model}} \times d_{ff}}$, $W_2 \in \R^{d_{ff} \times d_{\text{model}}}$, and typically $d_{ff} = 4 \cdot d_{\text{model}}$.
    \item \textbf{Add \& Norm} (second residual connection + normalization)
\end{enumerate}

\subsection{Decoder block}
\label{sec:decoder-block}

The decoder adds an additional layer:
\begin{enumerate}
    \item \textbf{Masked Multi-Head Self-Attention}\index{masked attention}: a causal mask prevents each position from ``observing'' future positions.
    \item \textbf{Add \& Norm}
    \item \textbf{Multi-Head Cross-Attention}\index{cross-attention}: the queries come from the decoder, the keys and values come from the encoder output.
    \item \textbf{Add \& Norm}
    \item \textbf{FFN}
    \item \textbf{Add \& Norm}
\end{enumerate}

\subsection{Pre-Norm vs Post-Norm}
\label{sec:pre-post-norm}

\begin{definition}[Post-Norm (original)]
\begin{equation}
x_{l+1} = \text{LayerNorm}(x_l + \text{SubLayer}(x_l))
\label{eq:post-norm}
\end{equation}
\end{definition}

\begin{definition}[Pre-Norm]
\begin{equation}
x_{l+1} = x_l + \text{SubLayer}(\text{LayerNorm}(x_l))
\label{eq:pre-norm}
\end{equation}
\end{definition}

\begin{proposition}[Gradient stability with Pre-Norm]
In the Pre-Norm case, the gradient flows directly through the residual connection without passing through the normalization:
\begin{equation}
\frac{\partial x_L}{\partial x_l} = I + \sum_{k=l}^{L-1} \frac{\partial \text{SubLayer}_k(\text{LayerNorm}(x_k))}{\partial x_l}
\end{equation}
The identity term $I$ guarantees a non-zero gradient flow, which stabilizes training for deep architectures ($L > 12$).
\end{proposition}

\begin{warning}
The original Transformer uses Post-Norm, but the majority of modern implementations (GPT-2, GPT-3, LLaMA) use Pre-Norm for its superior stability. A learning rate \emph{warm-up} is generally necessary with Post-Norm.
\end{warning}

\subsection{TikZ diagram of the complete architecture}

\begin{figure}[htbp]
\centering
\begin{tikzpicture}[
    block/.style={rectangle, draw, fill=blue!8, minimum width=3.2cm, minimum height=0.8cm, align=center, rounded corners=2pt},
    attnblock/.style={rectangle, draw, fill=orange!15, minimum width=3.2cm, minimum height=0.8cm, align=center, rounded corners=2pt},
    ffnblock/.style={rectangle, draw, fill=green!12, minimum width=3.2cm, minimum height=0.8cm, align=center, rounded corners=2pt},
    normblock/.style={rectangle, draw, fill=gray!15, minimum width=3.2cm, minimum height=0.6cm, align=center, rounded corners=2pt},
    io/.style={rectangle, draw, fill=yellow!15, minimum width=3.2cm, minimum height=0.7cm, align=center, rounded corners=2pt},
    arrow/.style={-{Stealth}, thick},
    dashed arrow/.style={-{Stealth}, thick, dashed},
    brace/.style={decorate, decoration={brace, amplitude=6pt, raise=4pt}},
]

% Encoder
\node[io] (emb-enc) at (0, 0) {Input Embedding\\+ Pos.\ Encoding};
\node[attnblock] (sa-enc) at (0, 1.5) {Multi-Head\\Self-Attention};
\node[normblock] (an1-enc) at (0, 2.6) {Add \& Norm};
\node[ffnblock] (ffn-enc) at (0, 3.6) {Feed-Forward};
\node[normblock] (an2-enc) at (0, 4.7) {Add \& Norm};

\draw[arrow] (emb-enc) -- (sa-enc);
\draw[arrow] (sa-enc) -- (an1-enc);
\draw[arrow] (an1-enc) -- (ffn-enc);
\draw[arrow] (ffn-enc) -- (an2-enc);

% Encoder brace
\draw[brace] (2.0, 1.1) -- (2.0, 5.1) node[midway, right=10pt, align=center] {$\times N$};

% Decoder
\node[io] (emb-dec) at (7, 0) {Output Embedding\\+ Pos.\ Encoding};
\node[attnblock] (msa-dec) at (7, 1.5) {Masked Multi-Head\\Self-Attention};
\node[normblock] (an1-dec) at (7, 2.6) {Add \& Norm};
\node[attnblock] (ca-dec) at (7, 3.6) {Multi-Head\\Cross-Attention};
\node[normblock] (an2-dec) at (7, 4.7) {Add \& Norm};
\node[ffnblock] (ffn-dec) at (7, 5.7) {Feed-Forward};
\node[normblock] (an3-dec) at (7, 6.8) {Add \& Norm};
\node[io] (out) at (7, 7.9) {Linear + Softmax};

\draw[arrow] (emb-dec) -- (msa-dec);
\draw[arrow] (msa-dec) -- (an1-dec);
\draw[arrow] (an1-dec) -- (ca-dec);
\draw[arrow] (ca-dec) -- (an2-dec);
\draw[arrow] (an2-dec) -- (ffn-dec);
\draw[arrow] (ffn-dec) -- (an3-dec);
\draw[arrow] (an3-dec) -- (out);

% Decoder brace
\draw[brace] (9.0, 1.1) -- (9.0, 7.2) node[midway, right=10pt, align=center] {$\times N$};

% Cross-attention arrow
\draw[dashed arrow] (an2-enc.east) -- ++(1.0, 0) |- (ca-dec.west)
    node[midway, above, font=\footnotesize] {$K, V$};

% Labels
\node[above=0.2cm, font=\bfseries] at (0, 5.1) {Encoder};
\node[above=0.2cm, font=\bfseries] at (7, 8.3) {Decoder};

\end{tikzpicture}
\caption{Complete Transformer architecture with encoder ($\times N$ blocks) and decoder ($\times N$ blocks). Dashed arrows represent the information flow from the encoder to the decoder's cross-attention mechanism.}
\label{fig:transformer-architecture}
\end{figure}

\section{Vision Transformer (ViT)}
\label{sec:vit}

The \emph{Vision Transformer}\index{Vision Transformer}\index{ViT} (Dosovitskiy et al., 2020) adapts the Transformer architecture to the computer vision domain.

\begin{definition}[Patch Embedding]
An image $x \in \R^{H \times W \times C}$ is split into $N = HW/P^2$ patches of size $P \times P$, each flattened and linearly projected:
\begin{equation}
z_0 = [\texttt{[CLS]};\; x_1^p E;\; x_2^p E;\; \ldots;\; x_N^p E] + E_{\text{pos}}
\end{equation}
where $E \in \R^{P^2 C \times d_{\text{model}}}$ is the projection matrix and $E_{\text{pos}} \in \R^{(N+1) \times d_{\text{model}}}$ is the learned positional encoding.
\end{definition}

\begin{remark}
The special \texttt{[CLS]}\index{CLS token} token serves as an aggregator: its final representation is used for classification.
\end{remark}

\section{BERT: Masked language modeling}
\label{sec:bert}

\textbf{BERT}\index{BERT} (Bidirectional Encoder Representations from Transformers, Devlin et al., 2019) uses only the \emph{encoder} of the Transformer.

\begin{definition}[MLM objective\index{Masked Language Modeling}]
15\% of the tokens in a sequence are randomly masked. The model must predict the masked tokens:
\begin{equation}
\mathcal{L}_{\text{MLM}} = -\E\!\left[\sum_{i \in \mathcal{M}} \log p(x_i \mid x_{\setminus \mathcal{M}};\, \theta)\right]
\end{equation}
where $\mathcal{M}$ is the set of masked positions.
\end{definition}

\section{GPT: Autoregressive language modeling}
\label{sec:gpt}

\textbf{GPT}\index{GPT} (Generative Pre-trained Transformer, Radford et al., 2018) uses only the \emph{decoder} with a causal mask\index{causal mask}.

\begin{definition}[Autoregressive objective]
The model predicts each token conditionally on the preceding tokens:
\begin{equation}
\mathcal{L}_{\text{AR}} = -\sum_{t=1}^{T} \log p(x_t \mid x_1, \ldots, x_{t-1};\, \theta)
\end{equation}
\end{definition}

The causal mask is implemented by adding $-\infty$ to future positions in the attention matrix before the $\softmax$:
\begin{equation}
\text{mask}_{ij} = \begin{cases} 0 & \text{if } j \leq i \\ -\infty & \text{if } j > i \end{cases}
\end{equation}

\section{PyTorch implementation of a Transformer}
\label{sec:transformer-pytorch}

\begin{codeblock}[title=\textbf{Complete Transformer for sequence classification}]
\begin{minted}{python}
import torch
import torch.nn as nn
import math

class PositionalEncoding(nn.Module):
    """Sinusoidal positional encoding."""
    def __init__(self, d_model, max_len=512, dropout=0.1):
        super().__init__()
        self.dropout = nn.Dropout(p=dropout)
        pe = torch.zeros(max_len, d_model)
        position = torch.arange(0, max_len).unsqueeze(1).float()
        div_term = torch.exp(
            torch.arange(0, d_model, 2).float()
            * (-math.log(10000.0) / d_model)
        )
        pe[:, 0::2] = torch.sin(position * div_term)
        pe[:, 1::2] = torch.cos(position * div_term)
        pe = pe.unsqueeze(0)  # (1, max_len, d_model)
        self.register_buffer('pe', pe)

    def forward(self, x):
        # x: (batch, seq_len, d_model)
        x = x + self.pe[:, :x.size(1)]
        return self.dropout(x)


class MultiHeadAttention(nn.Module):
    """Multi-head attention from scratch."""
    def __init__(self, d_model, n_heads):
        super().__init__()
        assert d_model % n_heads == 0
        self.d_k = d_model // n_heads
        self.n_heads = n_heads
        self.W_q = nn.Linear(d_model, d_model)
        self.W_k = nn.Linear(d_model, d_model)
        self.W_v = nn.Linear(d_model, d_model)
        self.W_o = nn.Linear(d_model, d_model)

    def forward(self, query, key, value, mask=None):
        B, T, D = query.shape
        # Projections and reshape: (B, n_heads, T, d_k)
        Q = self.W_q(query).view(B, T, self.n_heads, self.d_k).transpose(1, 2)
        K = self.W_k(key).view(B, -1, self.n_heads, self.d_k).transpose(1, 2)
        V = self.W_v(value).view(B, -1, self.n_heads, self.d_k).transpose(1, 2)

        # Scaled dot-product attention
        scores = torch.matmul(Q, K.transpose(-2, -1)) / math.sqrt(self.d_k)
        if mask is not None:
            scores = scores.masked_fill(mask == 0, float('-inf'))
        attn = torch.softmax(scores, dim=-1)
        out = torch.matmul(attn, V)

        # Concatenation and final projection
        out = out.transpose(1, 2).contiguous().view(B, T, D)
        return self.W_o(out)


class TransformerEncoderBlock(nn.Module):
    """Pre-Norm encoder block."""
    def __init__(self, d_model, n_heads, d_ff, dropout=0.1):
        super().__init__()
        self.attn = MultiHeadAttention(d_model, n_heads)
        self.ffn = nn.Sequential(
            nn.Linear(d_model, d_ff),
            nn.GELU(),
            nn.Dropout(dropout),
            nn.Linear(d_ff, d_model),
            nn.Dropout(dropout),
        )
        self.ln1 = nn.LayerNorm(d_model)
        self.ln2 = nn.LayerNorm(d_model)
        self.drop = nn.Dropout(dropout)

    def forward(self, x, mask=None):
        # Pre-Norm
        x_norm = self.ln1(x)
        x = x + self.drop(self.attn(x_norm, x_norm, x_norm, mask))
        x = x + self.ffn(self.ln2(x))
        return x


class TransformerClassifier(nn.Module):
    """Transformer encoder for sequence classification."""
    def __init__(self, vocab_size, d_model=128, n_heads=4,
                 d_ff=512, n_layers=4, n_classes=2,
                 max_len=256, dropout=0.1):
        super().__init__()
        self.embedding = nn.Embedding(vocab_size, d_model)
        self.pos_enc = PositionalEncoding(d_model, max_len, dropout)
        self.layers = nn.ModuleList([
            TransformerEncoderBlock(d_model, n_heads, d_ff, dropout)
            for _ in range(n_layers)
        ])
        self.ln_final = nn.LayerNorm(d_model)
        self.classifier = nn.Linear(d_model, n_classes)

    def forward(self, x, mask=None):
        # x: (batch, seq_len) integer indices
        x = self.embedding(x) * math.sqrt(self.embedding.embedding_dim)
        x = self.pos_enc(x)
        for layer in self.layers:
            x = layer(x, mask)
        x = self.ln_final(x)
        # Average over the sequence (alternative to [CLS] token)
        x = x.mean(dim=1)
        return self.classifier(x)


# --- Training ---
model = TransformerClassifier(
    vocab_size=10000, d_model=128, n_heads=4,
    d_ff=512, n_layers=4, n_classes=2
)
optimizer = torch.optim.AdamW(model.parameters(), lr=1e-4, weight_decay=0.01)
criterion = nn.CrossEntropyLoss()

# Example with synthetic data
batch_x = torch.randint(0, 10000, (32, 64))  # (batch=32, seq_len=64)
batch_y = torch.randint(0, 2, (32,))

logits = model(batch_x)
loss = criterion(logits, batch_y)
loss.backward()
optimizer.step()
print(f"Loss: {loss.item():.4f}, Logits shape: {logits.shape}")
\end{minted}
\end{codeblock}

\begin{codeoutput}
\begin{verbatim}
Loss: 0.6847, Logits shape: torch.Size([32, 2])
\end{verbatim}
\end{codeoutput}

\begin{bestpractice}
Always use \texttt{AdamW}\index{AdamW} with \emph{weight decay} for training Transformers. A \emph{learning rate schedule} with linear warm-up followed by cosine decay is standard.
\end{bestpractice}

\section{Ablation study}
\label{sec:transformer-ablation}

\begin{table}[htbp]
\centering
\caption{Ablation study on a Transformer for text classification (IMDB). Accuracy on the test set (\%).}
\label{tab:transformer-ablation}
\begin{tabular}{lcccc}
\toprule
\textbf{Configuration} & \textbf{Layers} & \textbf{Heads} & \textbf{$d_{ff}$} & \textbf{Accuracy (\%)} \\
\midrule
Base (Post-Norm) & 4 & 4 & 512 & 86.3 \\
Base (Pre-Norm) & 4 & 4 & 512 & 87.1 \\
\midrule
1 layer & 1 & 4 & 512 & 82.5 \\
2 layers & 2 & 4 & 512 & 85.4 \\
6 layers & 6 & 4 & 512 & 87.4 \\
8 layers & 8 & 4 & 512 & 87.2 \\
\midrule
1 head & 4 & 1 & 512 & 84.8 \\
2 heads & 4 & 2 & 512 & 86.0 \\
8 heads & 4 & 8 & 512 & 87.3 \\
\midrule
$d_{ff}=128$ & 4 & 4 & 128 & 84.1 \\
$d_{ff}=256$ & 4 & 4 & 256 & 85.9 \\
$d_{ff}=1024$ & 4 & 4 & 1024 & 87.5 \\
\bottomrule
\end{tabular}
\end{table}

\begin{remark}
The results show that: (1) Pre-Norm is systematically superior to Post-Norm without fine warm-up; (2) increasing the number of layers beyond 6 brings little benefit for this modestly sized dataset; (3) the FFN dimension has a significant impact; (4) a single attention head is clearly insufficient.
\end{remark}

\section{State of the art and connections}
\label{sec:transformers-sota}

\begin{sota}[title=\textbf{Large Language Models (LLMs)}]\index{LLM}\index{large language model}
Transformers are at the foundation of the LLM revolution:

\begin{itemize}
    \item \textbf{Scaling Laws}\index{scaling laws}: Kaplan et al.\ (2020) show that an LLM's performance follows a power law as a function of the number of parameters $N$, the dataset size $D$, and the compute budget $C$:
    \begin{equation}
    L(N) \approx \left(\frac{N_c}{N}\right)^{\alpha_N}, \quad \alpha_N \approx 0.076
    \end{equation}

    \item \textbf{GPT-4}\index{GPT-4} (OpenAI, 2023): presumably a Mixture of Experts (MoE)\index{Mixture of Experts} model with hundreds of billions of parameters.

    \item \textbf{Claude}\index{Claude} (Anthropic): a family of models using RLHF (Reinforcement Learning from Human Feedback)\index{RLHF} and Constitutional AI.

    \item \textbf{LLaMA}\index{LLaMA} (Meta, 2023--2024): open models using Pre-Norm with RMSNorm\index{RMSNorm}, RoPE\index{RoPE} (Rotary Position Embeddings), SwiGLU\index{SwiGLU} instead of ReLU, and Grouped Query Attention\index{GQA}.

    \item \textbf{Mixture of Experts (MoE)}\index{MoE}: Mixtral (Mistral, 2024) activates only 2 out of 8 experts per token, enabling a 47B parameter model with an inference cost of $\sim$13B.
\end{itemize}
\end{sota}

\begin{architecturebox}[title=\textbf{Evolution of Transformer architectures}]
\begin{center}
\begin{tikzpicture}[
    node distance=1.4cm,
    yearnode/.style={rectangle, draw, fill=blue!10, minimum width=2.5cm, minimum height=0.7cm, align=center, rounded corners=2pt, font=\small},
    arr/.style={-{Stealth}, thick},
]
\node[yearnode] (t0) {Transformer\\(2017)};
\node[yearnode, right=1.5cm of t0] (bert) {BERT\\(2018)};
\node[yearnode, above=0.8cm of bert] (gpt) {GPT / GPT-2\\(2018--2019)};
\node[yearnode, right=1.5cm of bert] (vit) {ViT\\(2020)};
\node[yearnode, right=1.5cm of gpt] (gpt3) {GPT-3\\(2020)};
\node[yearnode, right=1.5cm of vit] (llama) {LLaMA / Mistral\\(2023--2024)};
\node[yearnode, right=1.5cm of gpt3] (gpt4) {GPT-4 / Claude\\(2023--2025)};

\draw[arr] (t0) -- (bert);
\draw[arr] (t0) |- (gpt);
\draw[arr] (bert) -- (vit);
\draw[arr] (gpt) -- (gpt3);
\draw[arr] (gpt3) -- (gpt4);
\draw[arr] (vit) -- (llama);
\draw[arr] (gpt3) -- (llama);
\end{tikzpicture}
\end{center}
\end{architecturebox}

\section{Exercises}
\label{sec:transformers-exercises}

\begin{exercise}[Computing the complexity of attention]
\label{exo:attention-complexity}
\textbf{Difficulty: 1/3}
Show that the time complexity of scaled dot-product attention is $O(n^2 d_k)$ where $n$ is the sequence length and $d_k$ is the key dimension. Why does this pose a problem for long sequences? Cite two approaches to reduce this complexity.
\end{exercise}

\begin{exercise}[Causal mask]
\label{exo:causal-mask}
\textbf{Difficulty: 1/3}
Implement in PyTorch a function that generates a causal mask of size $n \times n$ and show its effect on the attention matrix for a sequence of 5 tokens.
\end{exercise}

\begin{exercise}[Number of parameters in a Transformer]
\label{exo:transformer-params}
\textbf{Difficulty: 2/3}
For a Transformer encoder with $L$ layers, $h$ heads, dimension $d_{\text{model}}$, and FFN dimension $d_{ff}$, compute the total number of parameters (not counting the embeddings).

\emph{Hint}: consider separately the multi-head attention ($4 d_{\text{model}}^2$), the FFN ($2 d_{\text{model}} d_{ff}$), and the LayerNorm ($4 d_{\text{model}}$) per block.
\end{exercise}

\begin{exercise}[ViT implementation]
\label{exo:vit-impl}
\textbf{Difficulty: 3/3}
Modify the code from Section~\ref{sec:transformer-pytorch} to create a Vision Transformer:
\begin{enumerate}
    \item Implement the \texttt{PatchEmbedding} that splits a $32 \times 32 \times 3$ image into patches of size $8 \times 8$.
    \item Add a learned \texttt{[CLS]} token.
    \item Train on CIFAR-10 and report the accuracy.
\end{enumerate}
\end{exercise}

\begin{exercise}[Comparison BERT vs GPT\@]
\label{exo:bert-vs-gpt}
\textbf{Difficulty: 2/3}
\begin{enumerate}
    \item Explain why BERT cannot be used directly for text generation.
    \item Show mathematically that BERT's MLM objective is not equivalent to maximizing $\log p(x)$ while GPT's autoregressive objective is.
    \item In which scenario would one prefer BERT over GPT and vice versa?
\end{enumerate}
\end{exercise}

\begin{exercise}[Analysis of attention heads {\normalfont (project)}]
\label{exo:attention-heads}
\textbf{Difficulty: 3/3}
Train a Transformer with 4 layers and 8 heads on a text classification task. Visualize the attention matrices of each head for example sentences. Can you identify specialized heads (position, syntax, semantics)? Evaluate the impact of pruning certain heads on performance.
\end{exercise}
